\section{As licenciaturas envolvidas na implementação da BNCC Computação}
\label{sec:diferentes-licenciaturas}

A BNCC Computação não estabelece que seus conhecimentos devam ser ensinados por um único perfil profissional. Também não exige que sejam organizados exclusivamente em uma disciplina autônoma. Dependendo da estratégia adotada por cada rede, sua implementação pode envolver professores especialistas, professores polivalentes e docentes de diferentes componentes curriculares.

Um levantamento sobre os currículos das unidades federativas identificou predomínio de estratégias transversais e uma presença ainda restrita de componentes obrigatórios específicos de Computação \cite{telefonica2025curriculos}. Em uma organização transversal, professores já responsáveis por outros componentes passam a participar da implementação. Em uma organização disciplinar, aumenta a necessidade de profissionais com formação mais específica.

O licenciado em Computação possui formação diretamente relacionada aos conhecimentos da área e pode atuar como professor especialista, formador ou colaborador de equipes curriculares. A adoção de um modelo baseado exclusivamente nesses profissionais, entretanto, é limitada pela escala de sua formação. Conforme discutido na Subseção~\ref{subsec:evidencias-contexto-brasileiro}, o número de cursos, concluintes e professores com formação específica em Computação é reduzido diante da abrangência da Educação Básica brasileira \cite{sbc2025grandesdesafios,sbc2022estatisticas,brasil2026emec}. Em uma análise dos microdados do Censo da Educação Superior de 2017 e 2018, \citeonline{menolli2021perfil} também identificaram uma redução de 36,3\% no número de ingressantes e um aumento da taxa de evasão de 16,02\% para 20,81\% nos cursos de Licenciatura em Computação.

Professores formados em Pedagogia atuam na Educação Infantil e nos Anos Iniciais, etapas nas quais a Computação pode ser desenvolvida por meio de atividades concretas, lúdicas e desplugadas. Seus conhecimentos sobre alfabetização, desenvolvimento infantil e organização do trabalho pedagógico constituem uma base importante para essa atuação. A formação nesses cursos, porém, geralmente não contempla de modo sistemático os conhecimentos específicos da Computação.

Professores de Matemática, Ciências e outras áreas podem participar de estratégias transversais ou interdisciplinares. Conceitos como lógica, modelagem, análise de dados e resolução de problemas oferecem possíveis pontos de articulação. Esses vínculos precisam ser tratados como pontos de entrada, e não como equivalentes aos conhecimentos próprios da Computação ou como garantia de cobertura de todos os seus eixos.

Esse cenário exige identificar quais perfis podem participar da implementação, quais responsabilidades podem assumir e que conhecimentos precisam desenvolver para desempenhá-las.

\subsection{Variação das necessidades formativas}
\label{subsec:variacao-necessidades-licenciaturas}

As necessidades formativas variam segundo a etapa da Educação Básica, o perfil de egresso da licenciatura, as possíveis articulações com a área de formação e a responsabilidade que o futuro professor deverá assumir.

Em Pedagogia, experiências com reconhecimento de padrões, resolução de problemas e jogos mostram que licenciandos podem desenvolver conhecimentos introdutórios de Pensamento Computacional sem formação técnica prévia \cite{correa2020pensamento}. Essas experiências fornecem evidências sobre a possibilidade de iniciar essa formação, enquanto a definição dos conteúdos, da profundidade e de sua distribuição no currículo permanece dependente dos objetivos estabelecidos para o curso.

Na Matemática, lógica, abstração, modelagem e algoritmos podem constituir pontos de articulação com a Computação. Antes de uma intervenção formativa, entretanto, a maioria dos licenciandos investigados por \citeonline{reichert2019pensamento} não reconhecia espontaneamente a relação entre Pensamento Computacional e conteúdos matemáticos. O resultado sugere que a proximidade entre as áreas precisa ser explicitada e desenvolvida durante a formação, em vez de ser tomada como um conhecimento já dominado pelos estudantes.


Nas Ciências, a Computação pode apoiar atividades de modelagem, simulação, coleta e análise de dados. Nas áreas de Linguagens e Ciências Humanas, as articulações podem envolver Cultura Digital, autoria, informação, participação social e impactos dos sistemas computacionais. Esses vínculos ajudam a contextualizar a formação, desde que a seleção curricular não restrinja a Computação apenas aos conhecimentos mais familiares a cada área.

Na Licenciatura em Computação, espera-se maior amplitude e profundidade conceitual, técnica e pedagógica. O curso pode oferecer referências sobre os conhecimentos próprios da área e formar especialistas capazes de colaborar com outras licenciaturas. Sua estrutura funciona, nesse caso, como fonte de conhecimentos e experiências, e não como modelo a ser reproduzido integralmente nos demais cursos.

As diferenças entre as licenciaturas envolvem mais do que a quantidade de conteúdo. Um curso pode preparar o egresso para introduzir aprendizagens de Computação, outro para relacioná-las a uma área disciplinar e outro para assumir o ensino especializado. Cada finalidade produz exigências próprias de conteúdo, profundidade, prática pedagógica e avaliação.

\subsection{Riscos de respostas curriculares uniformes}
\label{subsec:riscos-respostas-uniformes}

Uma resposta uniforme pode assumir duas formas. A primeira é considerar que somente licenciados em Computação devem ensinar os conhecimentos da área. Essa alternativa favorece maior especialização, mas encontra limites de escala e de distribuição dos profissionais. Conforme discutido na Subseção~\ref{subsec:evidencias-contexto-brasileiro}, a quantidade de professores com formação em Computação, de cursos ativos de Licenciatura em Computação e de concluintes formados anualmente é reduzida diante da abrangência da Educação Básica brasileira \cite{sbc2025grandesdesafios,sbc2022estatisticas,brasil2026emec}. Exigir um especialista em todas as escolas e etapas tornaria a implementação dependente de uma oferta profissional que não está disponível em escala compatível com o conjunto das redes de ensino.
A segunda forma é atribuir a responsabilidade a professores de outras áreas sem garantir preparação específica. Nesse caso, a Computação pode ser reduzida aos conteúdos já familiares ao docente ou ao uso de recursos digitais. Pensamento Computacional, Mundo Digital e Cultura Digital podem aparecer de forma fragmentada, sem progressão e sem cobertura suficiente. Uma revisão de 21 estudos sobre a preparação de professores dos anos iniciais mostrou que experiências formativas com participação ativa podem ampliar conhecimentos, atitudes e autoeficácia para ensinar Computação, indicando que essa preparação não deve ser presumida \cite{masonrich2019preparing}. Na Inglaterra, a implementação do currículo exigiu que professores generalistas dos anos iniciais desenvolvessem conhecimentos de conteúdo e de pedagogia da Computação e levou à criação de iniciativas nacionais de recrutamento, formação e apoio profissional \cite{fowlervegas2021england}. No Brasil, iniciativas de formação continuada em algoritmos e programação procuram preparar docentes de outras áreas diante da disponibilidade limitada de licenciados em Computação \cite{kretzer2020formacao}.

Modelos combinados procuram distribuir responsabilidades. Especialistas podem garantir maior densidade conceitual e apoiar outros docentes. Professores polivalentes e de diferentes áreas podem ampliar a cobertura e integrar conhecimentos aos contextos em que atuam. Entretanto, a combinação só é formativamente adequada quando as responsabilidades e os conhecimentos esperados de cada perfil são definidos. Experiências de coensino já relataram ganhos tanto para os alunos quanto para o desenvolvimento profissional dos próprios docentes envolvidos \cite{piech2020coteaching}, incluindo professores novatos de Computação que usam o coensino para preencher lacunas de conhecimento de conteúdo ao lado de colegas mais experientes \cite{yadav2024diving}.

Não é possível concluir, de forma geral, que uma única configuração seja adequada a todas as instituições. A existência de diferentes perfis reforça a necessidade de analisar cada licenciatura em relação a seu propósito, currículo e contexto.

\section{Dimensões da decisão curricular utilizadas na revisão}
\label{sec:decisao-curricular}

O Capítulo~\ref{chapter:referencial-teorico} definiu a integração curricular da Computação como um processo que relaciona os conhecimentos da área às finalidades, ao perfil profissional e às condições de uma licenciatura. Também mostrou que essa integração não se reduz à criação de uma disciplina ou à inclusão de conteúdos isolados, pois envolve decisões que se condicionam mutuamente.

Nesta seção, esses fundamentos são transformados em dimensões analíticas para examinar as respostas oferecidas pela literatura. As dimensões são utilizadas para identificar quais partes do processo cada contribuição ajuda a sustentar e quais permanecem pouco desenvolvidas.

A análise considera seis decisões principais: finalidade; conteúdos; profundidade; progressão; espaços curriculares; e práticas formativas e avaliação das aprendizagens. Também considera os atores responsáveis, as condições e restrições institucionais e a possibilidade de executar, institucionalizar e sustentar a alternativa proposta.

\subsection{Decisões envolvidas na integração curricular}
\label{subsec:decisoes-integracao-curricular}

As dimensões apresentadas a seguir podem ser examinadas separadamente para fins analíticos, mas precisam ser articuladas durante a elaboração de uma proposta curricular. Uma decisão altera as possibilidades das demais: a finalidade influencia os conhecimentos e a profundidade; os conteúdos e a progressão condicionam os espaços curriculares; e as práticas e formas de avaliação dependem dos objetivos estabelecidos.

A primeira envolve a \textbf{finalidade}. É preciso definir por que o futuro professor deverá aprender Computação e qual papel se espera que ele desempenhe. O objetivo pode ser compreender fundamentos da área, integrar conhecimentos de Computação a outro componente curricular, desenvolver atividades introdutórias ou assumir o ensino especializado. Essa definição orienta as decisões posteriores, pois uma formação voltada à familiarização com a área não exige a mesma amplitude, profundidade e experiência prática de uma formação destinada ao ensino especializado. Como os modelos internacionais variam de acordo com as políticas educacionais, os currículos escolares e as responsabilidades atribuídas aos diferentes perfis docentes, a finalidade não deve ser presumida nem transferida automaticamente de um contexto para outro \cite{yadav2022review}.

A segunda decisão trata dos \textbf{conteúdos}. A BNCC Computação oferece um referencial amplo, mas nem todos os cursos precisam abordar todos os conhecimentos com a mesma profundidade. A equipe deve identificar quais conhecimentos são pertinentes ao perfil do egresso e justificar a seleção realizada. Em comparação com currículos de outros países, a BNCC Computação brasileira optou por um nível de detalhamento conceitual mais elevado, com maior ênfase em aspectos sociais da área do que em aplicações práticas, o que já é, em si, uma decisão de conteúdo tomada em nível nacional \cite{eloy2024posicionando}. Em contraste, um levantamento dos Projetos Pedagógicos de licenciaturas em Ciências Biológicas de Pernambuco não encontrou nenhum curso que tivesse incorporado formalmente o Pensamento Computacional em sua estrutura curricular, mostrando que essa seleção de conteúdo às vezes simplesmente não chega a ser feita \cite{santos2022integracao}.

A terceira decisão diz respeito à \textbf{profundidade}. Um mesmo tema pode ser abordado como sensibilização, conhecimento introdutório, capacidade de aplicação pedagógica ou domínio especializado. A profundidade influencia a carga horária, os pré-requisitos e as atividades necessárias.

A quarta decisão envolve a \textbf{progressão}. Os conteúdos podem ser concentrados em um único momento ou distribuídos ao longo da formação. Uma organização progressiva permite relacionar fundamentos, práticas e experiências de estágio. Uma organização fragmentada pode produzir repetições ou lacunas.

A quinta decisão trata dos \textbf{espaços curriculares}. A Computação pode ser incluída em uma disciplina própria, distribuída entre componentes existentes, articulada às práticas e aos estágios ou organizada por uma combinação dessas alternativas. Um levantamento dos referenciais curriculares das 27 unidades federativas mostrou que, ao final de 2024, apenas quatro possuíam componente obrigatório relacionado à Computação e somente duas dessas quatro contemplavam os três eixos da BNCC Computação. Embora a transversalidade fosse mencionada em todos os referenciais estaduais, na maioria deles as referências permaneciam restritas à competência geral de Cultura Digital e a temas transversais, sem incorporar habilidades e conteúdos estruturantes da área \cite{telefonica2025curriculos}.

Uma disciplina própria favorece a identificação do conteúdo, a definição de responsabilidade e a avaliação de sua oferta. Entretanto, exige espaço na carga horária e profissionais capazes de ministrá-la. Uma inserção transversal pode aproveitar componentes existentes, mas exige coordenação para evitar menções isoladas e sem continuidade.

A sexta decisão envolve as \textbf{práticas formativas e a avaliação das aprendizagens}. Uma formação voltada ao ensino não pode se restringir à apresentação de conceitos. Conforme a finalidade estabelecida, os licenciandos precisam de oportunidades para resolver problemas, analisar atividades, elaborar planos de aula, experimentar estratégias didáticas e refletir sobre sua aplicação. A revisão de \citeonline{rodrigues2024integration} indica a necessidade de formações mais prolongadas e integradas ao currículo, que articulem componentes teóricos, experiências práticas e reflexão sobre a prática docente.

\subsection{Atores, condições e restrições}
\label{subsec:atores-condicoes-restricoes}

As decisões curriculares são tomadas por atores situados em instituições. Segundo a \sigla{CONAES}{Comissão Nacional de Avaliação da Educação Superior} de \citeyear{brasil2010resolucaoconaes}, o \sigla{NDE}{Núcleo Docente Estruturante}, constituído por docentes do curso, atua no processo de concepção, consolidação e contínua atualização do Projeto Pedagógico de Curso. A elaboração, a revisão e a aprovação de mudanças curriculares seguem a organização de cada instituição e podem envolver, além do NDE, a coordenação e o colegiado do curso, conselhos departamentais e o conselho superior da instituição \cite{silva2020gestao}.

Esse processo envolve negociação. Inserir novos conhecimentos pode exigir a criação de componentes, a revisão de ementas, o deslocamento de cargas horárias ou a redistribuição de responsabilidades. Como o currículo possui limites, a inclusão de uma demanda pode afetar espaços já ocupados por outras áreas. Um estudo de caso numa universidade comunitária catarinense mostra como esse processo pode se desenrolar em grande escala: a revisão de 64 projetos pedagógicos de curso envolveu Núcleos Docentes Estruturantes, colegiados e uma comissão vinculada à Pró-Reitoria de Ensino, resultando na criação de componentes curriculares comuns a diferentes cursos e na redistribuição da carga horária entre eles \cite{krahe2013reforma}. Um estudo com 41 docentes vinculados a 13 departamentos de faculdades de Educação na Turquia ilustra algumas dificuldades desse processo. Os participantes relataram tanto resistência quanto iniciativas de mudança construídas de baixo para cima, geralmente concentradas em ajustes de componentes individuais. A maioria também reconheceu dificuldades e falta de preparação profissional para desenvolver e avaliar currículos \cite{aydinok2025behind}.

As alternativas também são condicionadas por normas. Diretrizes curriculares, carga horária total, organização dos núcleos formativos e regras internas delimitam o que pode ser alterado e quais procedimentos precisam ser seguidos.

Há ainda condições humanas. A instituição pode não dispor de docentes com formação em Computação ou com preparação pedagógica para ensinar seus conhecimentos em uma licenciatura. Contratar novos profissionais, formar o corpo docente existente ou estabelecer colaboração entre departamentos e unidades acadêmicas produzem custos, prazos e possibilidades diferentes.

Uma condição específica diz respeito à preparação dos docentes responsáveis por formar os licenciandos. Professores da área de Computação podem dominar os conhecimentos específicos, mas não conhecer suficientemente as finalidades, as práticas pedagógicas e o perfil formativo de outra licenciatura. Em sentido inverso, docentes da área educacional ou dos componentes de métodos de ensino podem conhecer o curso e seus estudantes, mas não possuir formação suficiente em Computação. No contexto brasileiro, \citeonline{falcaofranca2021computational} defendem tanto o desenvolvimento dos conhecimentos de Computação dos formadores quanto a colaboração entre docentes das licenciaturas em Computação e de outras áreas. \citeonline{yadav2021building} também destacam que a incorporação da Computação à formação inicial depende da ampliação da capacidade das instituições formadoras, por meio da preparação do corpo docente e de parcerias entre áreas de Educação, Computação e redes escolares. Assim, a análise curricular precisa identificar quem poderá assumir cada responsabilidade, quais conhecimentos esses docentes já possuem, que preparação adicional será necessária e como a colaboração poderá ser mantida ao longo do tempo.

As condições materiais também afetam as alternativas. Determinados objetivos podem exigir laboratórios, dispositivos, conectividade, programas ou apoio técnico. Outros podem ser alcançados por meio de atividades desplugadas ou com recursos já disponíveis. A análise precisa distinguir requisitos indispensáveis de recursos desejáveis.

\subsection{Viabilidade, institucionalização e sustentabilidade}
\label{subsec:viabilidade-institucionalizacao-sustentabilidade}

A distinção entre adoção formal e institucionalização, discutida na Subseção~\ref{subsec:desafios-interdependentes}, permite examinar se uma proposta curricular apresenta apenas uma configuração possível ou se também considera as condições necessárias para que ela seja realizada e incorporada ao funcionamento do curso.

Uma disciplina pode constar no Projeto Pedagógico de Curso e não ser oferecida por falta de docente. Um conteúdo pode aparecer em uma ementa e não ser tratado durante as aulas. Uma atividade pode ocorrer durante um projeto específico e desaparecer quando o financiamento, a equipe ou as condições institucionais se modificam. Essa distância entre alteração formal e mudança nas práticas foi observada em um estudo com 17 docentes da Universidade do Porto após a adoção do Processo de Bolonha. Os participantes reconheceram transformações importantes na organização dos cursos e das unidades curriculares, enquanto consideraram menores as mudanças nos aspectos operacionais do exercício da docência, do ensino-aprendizagem e da avaliação \cite{leiteramos2015reconfiguracoes}.

Para analisar as contribuições da literatura, esta pesquisa distingue três aspectos. A \textit{viabilidade} diz respeito à possibilidade de executar uma alternativa diante das normas, da carga horária, dos docentes, dos recursos e dos procedimentos institucionais. A \textit{institucionalização} envolve sua incorporação aos documentos, às responsabilidades e às práticas regulares do curso. A \textit{sustentabilidade} se refere às condições para que a mudança permaneça quando os participantes, os recursos ou as circunstâncias iniciais se modificam.

Uma proposta pode ser pedagogicamente adequada e, ao mesmo tempo, inviável nas condições atuais de um curso. Também pode ser viável como experiência pontual sem possuir condições de institucionalização ou continuidade. Essas diferenças precisam ser explicitadas para que a equipe curricular reconheça não apenas a alternativa desejada, mas também suas dependências, seus riscos e as mudanças intermediárias que podem ser necessárias.

Integrar Computação a uma licenciatura constitui, portanto, um problema de decisão contextualizada. A equipe precisa relacionar finalidades, conteúdos, profundidade, progressão, espaços curriculares, práticas formativas e avaliação às condições reais do curso. A seção seguinte examina em que medida os currículos atuais e as contribuições da literatura já apoiam esse processo.

\section{A presença da Computação nos currículos e as respostas oferecidas pela literatura}
\label{sec:situacao-solucoes}

A literatura examinada nesta seção foi reunida por meio de uma revisão narrativa orientada pelo recorte temático da pesquisa \cite{Ogassavara2025}. Os trabalhos foram localizados por buscas livres e pelo encadeamento de referências a partir de publicações consideradas pertinentes. A seleção priorizou estudos que contribuíssem diretamente para compreender a presença da Computação nos currículos de licenciatura, os conhecimentos necessários à formação docente, as possibilidades de integração curricular ou os processos institucionais envolvidos na mudança dos cursos. Não foi estabelecido um período de publicação nem adotado um protocolo sistemático de busca, triagem e exclusão. Assim, o conjunto analisado não pretende ser exaustivo, mas reunir contribuições relevantes para caracterizar o problema, comparar as respostas existentes e fundamentar a lacuna investigada.

Estudos documentais brasileiros indicam que a presença da Computação nos currículos de licenciatura é heterogênea. Os resultados variam entre cursos, instituições e áreas, mas apontam problemas recorrentes de ausência, optatividade, baixa carga horária, falta de progressão e associação com o uso instrumental de tecnologias.

\citeonline{silvafalcao2021pesquisa} analisaram Projetos Pedagógicos de Licenciaturas em Computação e identificaram um processo de transformação curricular ainda em andamento. Mesmo nesses cursos, diretamente vinculados à área, existem tensões entre formação técnica, formação pedagógica e definição do perfil profissional.

Em cursos de Ciências Biológicas de Pernambuco, \citeonline{santos2022integracao} não identificaram a incorporação formal do Pensamento Computacional nos Projetos Pedagógicos analisados. O resultado mostra que a presença de tecnologias educacionais não garante a inclusão de conhecimentos próprios da Computação.

\citeonline{abrantes2024pensamento} analisaram cursos de Licenciatura em Matemática de instituições federais do Nordeste. O estudo identificou componentes relacionados à informática e às tecnologias, mas apontou a necessidade de maior clareza curricular e de reformulação para incorporar o Pensamento Computacional de forma adequada.

Esses estudos cumprem uma função diagnóstica. Eles mostram se a Computação aparece nos currículos, em quais componentes, com que obrigatoriedade e por meio de quais conteúdos. Também permitem identificar situações em que componentes denominados ``Informática em Educação'' ou ``Tecnologias da Informação e Comunicação'' permanecem centrados no uso de ferramentas, sem contemplar algoritmos, sistemas, dados ou práticas de Pensamento Computacional.

\subsection{Revisões e referenciais que caracterizam o campo}
\label{subsec:revisoes-referenciais-campo}

A literatura internacional sobre Computação na formação inicial reúne revisões, referenciais de conteúdo e frameworks voltados à formação docente.

\citeonline{listiaji2025integrating} revisaram 43 artigos publicados entre 2014 e 2024 e identificaram concentração em áreas \sigla{STEM}{\textit{Science, Technology, Engineering and Mathematics}}, atividades plugadas, programação visual e avaliações baseadas em questionários e testes. O resultado mostra que existe uma produção consistente sobre intervenções formativas, mas também evidencia concentração temática e disciplinar.

\citeonline{rodrigues2024integration}, em uma revisão sobre a integração do Pensamento Computacional na formação inicial de professores dos anos iniciais, identificaram predomínio de módulos inseridos em componentes existentes e de cursos optativos, muitos deles de curta duração. Os autores defenderam uma formação prolongada, integrada ao currículo e articulada ao planejamento, à prática e à reflexão docente. \citeonline{yun2025computational}, ao revisarem pesquisas sobre a integração do Pensamento Computacional à formação inicial de professores de Ciências, caracterizaram o campo como predominantemente exploratório, com forte presença de desenhos observacionais, amostras geralmente pequenas e concentração em determinadas áreas das Ciências, sobretudo nas Ciências Físicas.

\citeonline{yadav2022review} analisaram modelos internacionais de formação de professores de Computação e mostraram que as configurações variam conforme políticas, currículos escolares e condições de cada país. Essa variação reforça que modelos não podem ser transferidos sem considerar o contexto institucional e normativo.

Frameworks também ajudam a identificar conhecimentos, capacidades e aspectos que precisam ser considerados na formação docente. \citeonline{angeli2016k6} propuseram um framework curricular de Pensamento Computacional para a educação K--6 e discutiram suas implicações para os conhecimentos necessários aos professores. O EnCITE, por sua vez, organiza os problemas e as oportunidades envolvidos na integração da Computação e das tecnologias digitais à formação de professores em cinco espaços interligados: tecnológico, pedagógico, definicional e ideológico, político e de desenvolvimento \cite{vogel2024encite}.

Essas contribuições delimitam conhecimentos docentes, conteúdos possíveis e aspectos que precisam ser observados na integração. Elas ajudam a responder o que pode ser ensinado, que capacidades o professor precisa desenvolver e que tipos de problema podem interferir na formação. Sua função predominante é caracterizar os elementos relevantes para o planejamento. A passagem desses elementos para uma proposta de alteração de um currículo específico demanda outro tipo de procedimento, capaz de relacioná-los ao perfil e às condições do curso.

\subsection{Intervenções e configurações curriculares}
\label{subsec:intervencoes-configuracoes-curriculares}

Outro conjunto de contribuições reúne intervenções implementadas, propostas de disciplinas, sequências formativas e modelos de desenvolvimento institucional.

\citeonline{yadav2014computational} avaliaram a inserção de módulos introdutórios de Pensamento Computacional em uma disciplina obrigatória da formação inicial de professores. \citeonline{yadav2017computational} recomendaram o redesenho de disciplinas de tecnologia educacional e a utilização de disciplinas de métodos de ensino como espaços para desenvolver e relacionar o Pensamento Computacional às diferentes áreas de conteúdo. Esses trabalhos mostram que conhecimentos de Computação podem ser incorporados a espaços curriculares existentes, seja por meio de intervenções concretas, seja pela formulação de pontos de entrada para a formação docente.

Também integram esse conjunto duas experiências brasileiras discutidas anteriormente na Subseção~\ref{subsec:variacao-necessidades-licenciaturas}. \citeonline{correa2020pensamento} desenvolveram uma experiência formativa com licenciandos de Pedagogia, enquanto \citeonline{reichert2019pensamento} investigaram uma intervenção com licenciandos de Matemática. Embora não correspondam a reformas curriculares completas, esses trabalhos implementam atividades formativas com futuros professores e oferecem evidências sobre possibilidades e limites da integração em licenciaturas com perfis distintos.

A coletânea organizada por \citeonline{mouza2021preparing} reúne modelos de cursos, práticas formativas, propostas institucionais e discussões sobre políticas de formação de professores de Ciência da Computação. Seus capítulos apresentam contribuições de naturezas distintas e não constituem, em conjunto, uma única intervenção.

\citeonline{lucarelli2021creating} descrevem um modelo de formação destinado a preparar futuros professores de diferentes áreas para facilitar conceitos de Ciência da Computação. \citeonline{yang2021redesigning} analisam o redesenho de uma disciplina de tecnologia educacional para integrar o Pensamento Computacional ao ensino das áreas de conteúdo. \citeonline{odombartel2021preparing} apresentam um modelo de dois cursos para licenciandos em Matemática relacionado ao currículo de \textit{AP Computer Science Principles}. \citeonline{yadav2021building} discutem ações por meio das quais as escolas de Educação podem ampliar sua capacidade institucional de preparar e apoiar professores de Ciência da Computação.

Em conjunto, esses trabalhos documentam diferentes possibilidades de inserção da Computação na formação inicial. Alguns apresentam componentes e experiências implementadas; outros propõem modelos de curso, pontos de entrada ou estratégias institucionais. Eles mostram que a integração pode ocorrer em disciplinas obrigatórias, componentes de tecnologia educacional, disciplinas de métodos, sequências específicas ou iniciativas distribuídas ao longo da formação.

Seus limites decorrem, em parte, de sua especificidade e de seus diferentes estágios de desenvolvimento e avaliação. Uma sequência construída para licenciandos em Matemática ou relacionada a um currículo norte-americano responde a finalidades e condições próprias. Ela pode orientar outras instituições, mas não determina como deve ser adaptada a diferentes licenciaturas brasileiras, às exigências da BNCC Computação ou aos recursos disponíveis em cada curso.

\subsection{Propostas e processos mais próximos desta pesquisa}
\label{subsec}

Alguns trabalhos se aproximam mais diretamente do problema de apoiar decisões de integração curricular, seja por investigarem a colaboração entre áreas, seja por proporem conteúdos, configurações ou processos de mudança.

\citeonline{santosfalcao2024pensamento} investigaram oportunidades e dificuldades percebidas por docentes de diferentes licenciaturas. A partir dos resultados, recomendaram a colaboração entre docentes de Licenciatura em Computação e integrantes dos Núcleos Docentes Estruturantes de outros cursos para analisar Projetos Pedagógicos e conceber formas de integração adequadas a cada área. A contribuição do estudo está em reconhecer que a integração precisa ser construída conjuntamente por participantes que conheçam tanto a Computação quanto o currículo da licenciatura.

\citeonline{mendes2024proposta} propõe recomendações para incluir Pensamento Computacional nos currículos de licenciatura. A publicação descreve uma pesquisa em andamento, fundamentada em revisão da literatura e com previsão de adaptação das recomendações a um currículo real. Seu foco está na organização de orientações curriculares relacionadas ao Pensamento Computacional.

\citeonline{delgado2026computacao} propõe um arcabouço curricular mínimo de Computação para a formação inicial de professores. A proposta distribui conhecimentos entre componentes obrigatórios, formação específica, práticas pedagógicas e componentes optativos. Seu escopo contempla Pensamento Computacional, Mundo Digital, Cultura Digital e Inteligência Artificial, além de reconhecer que a configuração precisa variar entre as licenciaturas.

\citeonline{margulieux2022integrating} descrevem um processo de \textit{design-based research} para integrar Computação a programas de formação de professores de Linguagens, Matemática e Ciências. O processo reuniu especialistas em Computação e docentes das áreas para definir objetivos, selecionar ou desenvolver atividades, inseri-las em disciplinas de métodos de ensino e avaliar as experiências com os licenciandos. Entre os trabalhos examinados, essa experiência oferece o exemplo empírico internacional mais próximo de um processo colaborativo de transformação curricular que considera diferenças entre áreas.

Os trabalhos se aproximam do problema desta pesquisa por caminhos complementares. \citeonline{santosfalcao2024pensamento} e \citeonline{margulieux2022integrating} evidenciam a colaboração necessária para analisar e modificar os currículos. \citeonline{mendes2024proposta} e \citeonline{delgado2026computacao} concentram-se na definição de conhecimentos, recomendações e possíveis configurações formativas. O primeiro grupo esclarece quem precisa participar e como a mudança pode ser construída; o segundo oferece referências sobre o que pode compor a formação.

A principal lacuna aparece na ligação entre esses dois conjuntos de contribuições. \citeonline{santosfalcao2024pensamento} recomendam a colaboração, sem desenvolver e avaliar o instrumento que a conduziria. O processo documentado por \citeonline{margulieux2022integrating} foi realizado em uma instituição e concentrou-se principalmente em programação e Pensamento Computacional. As propostas de \citeonline{mendes2024proposta} e \citeonline{delgado2026computacao} apresentam orientações ou configurações curriculares, mas não oferecem um procedimento empiricamente avaliado para que diferentes equipes relacionem esses referenciais ao currículo vigente, comparem alternativas e justifiquem como conteúdos, profundidade, espaços e progressão devem variar em cada contexto.

\subsection{Síntese comparativa das contribuições}
\label{subsec:sintese-comparativa-literatura}

Os trabalhos reunidos por esse procedimento oferecem contribuições de naturezas distintas para a integração da Computação na formação inicial de professores. Para fins desta pesquisa, essas contribuições foram agrupadas analiticamente segundo a função predominante que desempenham diante do problema investigado.

O agrupamento constitui uma síntese temática elaborada pelo autor para apoiar a comparação dos trabalhos. Ele não corresponde a uma tipologia consolidada na literatura nem resulta de um procedimento formal de codificação.

A classificação considera como unidade de análise cada contribuição individual. No caso de coletâneas, os capítulos são classificados separadamente, pois uma mesma obra pode reunir intervenções implementadas, modelos de curso e propostas de desenvolvimento institucional. Quando um trabalho apresenta características de mais de uma categoria, sua classificação principal corresponde ao produto ou à função que predomina na análise desenvolvida nesta seção.

A coletânea organizada por \citeonline{mouza2021preparing} exemplifica essa heterogeneidade ao reunir modelos de formação, experiências de redesenho de disciplinas, sequências de cursos e estratégias de ampliação da capacidade institucional. Por essa razão, seus capítulos são considerados conforme a natureza específica de cada contribuição.

Para compor a síntese, foram considerados os trabalhos que oferecem uma contribuição direta para compreender, planejar ou realizar a integração da Computação em currículos de formação inicial de professores. Isso inclui estudos que analisam currículos de licenciatura, organizam conhecimentos ou dimensões relevantes, implementam experiências com licenciandos, propõem configurações curriculares ou examinam processos institucionais de mudança.

A Tabela ~\ref{quad:comparacao-contribuicoes} apresenta os seis tipos de contribuição, o que cada um oferece e os principais trabalhos que atendem a esse recorte.


\begin{table}[htbp]
  \centering
  \caption{Síntese das contribuições relacionadas à integração da Computação em licenciaturas}
  \label{quad:comparacao-contribuicoes}
  \small
  \setlength{\tabcolsep}{4pt}
  \renewcommand{\arraystretch}{1.2}
  \begin{tabularx}{\textwidth}{
    >{\raggedright\arraybackslash}p{3.0cm}
    >{\raggedright\arraybackslash}X
    >{\raggedright\arraybackslash}p{5.5cm}
  }
    \toprule
    \textbf{Tipo de contribuição} & \textbf{O que oferece} & \textbf{Estudos relacionados} \\
    \midrule

    Diagnósticos curriculares e análises documentais &
    Caracterizam a presença da Computação em \sigla{PPCs}{Projeto Pedagógico de Curso}, matrizes e ementas, mostrando ausências, presença parcial e diferenças entre cursos. &
    \cite{silvafalcao2021pesquisa,santos2022integracao,abrantes2024pensamento} \\

    \addlinespace

    Revisões e sínteses do campo &
    Reúnem resultados de vários estudos sobre formação inicial em Computação, revelando padrões e lacunas do campo. &
    \cite{listiaji2025integrating,rodrigues2024integration,yun2025computational,yadav2022review} \\

    \addlinespace

    Frameworks e referenciais analíticos &
    Organizam conhecimentos ou espaços de problemas e oportunidades que precisam ser considerados na integração da Computação à formação docente. & \cite{angeli2016k6,vogel2024encite} \\

    \addlinespace

    Intervenções e experiências formativas &
    Mostram atividades, módulos ou disciplinas já aplicados com licenciandos de diferentes áreas. &
    \cite{yadav2014computational,yang2021redesigning,correa2020pensamento,reichert2019pensamento} \\

    \addlinespace

    Modelos e propostas curriculares &
    Propõem conteúdos, componentes e sequências para organizar a Computação ao longo da formação docente. &
    \cite{yadav2017computational,lucarelli2021creating,odombartel2021preparing,mendes2024proposta,delgado2026computacao} \\

    \addlinespace

    Processos de mudança, colaboração e institucionalização curricular &
    Descrevem como especialistas, docentes e instâncias institucionais negociam e sustentam mudanças curriculares. &
    \cite{margulieux2022integrating,santosfalcao2024pensamento,yadav2021building,silva2020gestao,krahe2013reforma,kellerfranco2018formacao} \\

    \bottomrule
  \end{tabularx}
  \fonte{Elaborado pelo autor.}
\end{table}

Vistas em conjunto, as seis categorias cumprem funções complementares diante do problema investigado. Os diagnósticos e as revisões caracterizam o ponto de partida; os frameworks e modelos delimitam conhecimentos, dimensões e possíveis configurações; as intervenções mostram formas de inserção já experimentadas; e os estudos sobre colaboração e mudança curricular esclarecem como diferentes participantes podem construir e sustentar alterações nos cursos.

Os diagnósticos curriculares mostram o que já está presente ou ausente. \citeonline{santos2022integracao}, por exemplo, não encontraram o Pensamento Computacional formalmente incorporado nos Projetos Pedagógicos de licenciaturas em Ciências Biológicas de Pernambuco. As revisões ampliam essa caracterização ao revelar padrões do campo, como a concentração de estudos em áreas STEM e o predomínio de ações formativas de curta duração \cite{listiaji2025integrating}. Essas contribuições estabelecem o problema que precisa ser enfrentado e fornecem referências para avaliar a situação de um curso.

Frameworks e modelos curriculares oferecem insumos para formular alternativas. O EnCITE reúne problemas e oportunidades de natureza tecnológica, pedagógica, definicional e ideológica, política e de desenvolvimento \cite{vogel2024encite}. \citeonline{delgado2026computacao} apresenta uma possível distribuição de componentes que contempla os três eixos da BNCC Computação e também a Inteligência Artificial. O valor dessas contribuições está na explicitação de aspectos, conhecimentos e configurações que uma equipe pode considerar durante o planejamento.

As intervenções formativas documentam possibilidades de aplicação com licenciandos. \citeonline{correa2020pensamento} descrevem uma experiência com estudantes de Pedagogia, enquanto \citeonline{reichert2019pensamento} examinam uma intervenção com estudantes de Matemática. Esses trabalhos fornecem evidências sobre atividades, dificuldades e aprendizagens que podem ocorrer em cursos nos quais a Computação não constitui a área principal de formação.

Os estudos sobre mudança e colaboração curricular tratam do processo institucional. \citeonline{margulieux2022integrating} mostram especialistas em Computação e docentes de outras áreas definindo objetivos, selecionando atividades e inserindo-as em disciplinas de métodos de ensino. \citeonline{santosfalcao2024pensamento} indicam a aproximação entre docentes de Computação e integrantes de Núcleos Docentes Estruturantes como caminho para analisar os Projetos Pedagógicos de diferentes licenciaturas.

A comparação revela uma cobertura desigual das decisões necessárias. A literatura oferece informações consistentes sobre conhecimentos possíveis, experiências já realizadas, espaços de inserção e participantes envolvidos. São menos desenvolvidos os critérios para relacionar essas informações ao perfil de um curso, definir a profundidade e a progressão, comparar configurações curriculares e examinar sua viabilidade diante das condições institucionais.

Uma equipe curricular pode reunir a BNCC Computação, o currículo vigente, frameworks, propostas formativas e relatos de experiências. O ponto ainda pouco estruturado é o processo de transformar esse conjunto em alternativas contextualizadas, estabelecer critérios de comparação e registrar as razões pelas quais determinada configuração foi escolhida. É essa insuficiência que a seção seguinte formaliza.

\section{A lacuna que fundamenta a pesquisa}
\label{sec:lacuna}

A Seção~\ref{sec:decisao-curricular} mostrou que integrar a Computação a uma licenciatura não corresponde apenas a selecionar conteúdos. Esse processo envolve decisões interdependentes sobre as finalidades da formação, os conhecimentos a serem contemplados, sua profundidade e progressão, os espaços curriculares em que serão inseridos, as práticas formativas adotadas e as condições institucionais necessárias para viabilizá-los. A Seção~\ref{sec:situacao-solucoes}, por sua vez, mostrou que a literatura oferece contribuições para diferentes partes desse processo, incluindo diagnósticos curriculares, revisões, frameworks analíticos, intervenções formativas, propostas curriculares e experiências de mudança colaborativa.

O que permanece insuficientemente desenvolvido, no conjunto de trabalhos examinado, é um processo explícito, reutilizável e avaliado que apoie a passagem dessas contribuições gerais para decisões curriculares contextualizadas em licenciaturas com diferentes finalidades, estruturas e condições de funcionamento.

As contribuições existentes ajudam a responder partes do problema. Os diagnósticos permitem reconhecer o que já está presente no currículo; os referenciais delimitam conhecimentos que podem ser considerados; as intervenções e os modelos curriculares apresentam possibilidades de inserção; e os estudos sobre mudança curricular evidenciam a importância da colaboração, da negociação e da sustentação institucional. Esses elementos, entretanto, não aparecem articulados em um procedimento que oriente uma equipe curricular a reunir as informações necessárias, relacioná-las ao perfil do curso, formular alternativas, examinar sua viabilidade e justificar as escolhas realizadas.

Permanece pouco estruturado, em particular, o processo pelo qual demandas normativas e formativas são transformadas em decisões sobre o que deve ser integrado, com qual finalidade, em que profundidade, em qual momento do curso, por meio de quais componentes ou práticas e sob quais condições. Também não está suficientemente operacionalizada a comparação entre alternativas, como a criação de um componente específico, a inserção em disciplinas existentes, a distribuição dos conhecimentos ao longo do curso ou a combinação dessas possibilidades.

Essa insuficiência possui consequências práticas. Quando o processo de análise não é explicitado, as decisões podem permanecer dependentes do conhecimento tácito e da experiência individual dos participantes. A equipe pode adotar uma proposta incompatível com a estrutura do curso, restringir a Computação aos conteúdos mais próximos das áreas já presentes no currículo ou formular recomendações que desconsiderem a disponibilidade de docentes, a carga horária, os recursos materiais e os procedimentos institucionais de aprovação. A ausência de registros estruturados também dificulta compreender por que determinada alternativa foi escolhida, de quais condições sua execução depende e em que circunstâncias ela precisaria ser revista.

A lacuna situa-se na mediação entre o conhecimento disponível e a decisão curricular local. A proposta desta pesquisa se insere, portanto, no espaço de investigar como essa mediação pode ser estruturada de modo que equipes curriculares produzam recomendações diferenciadas, justificadas e condicionadas às características de cada curso.

\section{Considerações finais}
\label{sec:consideracoes-finais-revisao-integracao}

Este capítulo mostrou que a integração da Computação nas licenciaturas precisa variar conforme o perfil do curso, as responsabilidades esperadas de seus egressos e suas condições institucionais. Por isso, não pode ser reduzida à adoção de uma disciplina ou de um conjunto uniforme de conteúdos. Ela envolve decisões articuladas sobre finalidade, conhecimentos, profundidade, progressão, espaços curriculares, práticas formativas e avaliação.

A literatura examinada oferece diagnósticos, referenciais, intervenções, modelos curriculares e experiências de mudança institucional. Essas contribuições apoiam partes do processo, mas não apresentam, de forma articulada, um procedimento que ajude equipes curriculares a relacioná-las ao currículo vigente, comparar alternativas e justificar decisões adequadas às condições de cada licenciatura.

Essa insuficiência fundamenta a proposta apresentada no capítulo seguinte: o desenvolvimento e a avaliação de um protocolo sistematizado de apoio à decisão curricular, implementado em uma ferramenta computacional, para orientar a integração contextualizada da Computação em cursos de licenciatura.